% ============================================================
%  CHAPTER 1 — Introduction
% ============================================================
\chapter{Introduction}
\label{ch:introduction}

\section{Context and Motivation}

We are living through an inflection point in artificial intelligence.
The 2017 Transformer architecture \cite{vaswani2017attention} sparked a revolution in
scale: models trained on internet-scale corpora --- GPT-3 \cite{brown2020gpt3},
GPT-4 \cite{achiam2023gpt4}, Gemini \cite{team2023gemini} --- demonstrate reasoning,
planning, and instruction-following capabilities that were unimaginable five years prior.
The discovery of \emph{emergent abilities} --- qualitatively new skills that appear
suddenly at sufficient scale, absent in smaller models --- suggests that these systems
are not merely doing sophisticated pattern matching but developing internal world models.

When vision encoders were fused with these language models to produce
\textbf{Vision-Language Models} (VLMs), the implications for robotics became
immediately clear.
A VLM can look at a kitchen scene and say ``there is a bottle to the left of the sink,
and the person is standing near the window'' --- precisely the spatial understanding
that a robot needs to act purposefully.
The \textbf{agentic paradigm} completed the picture: by structuring VLM inference as a
loop of perception, chain-of-thought reasoning, skill selection, execution, and
observation (the ReAct pattern \cite{yao2022react}), a VLM becomes not just an
answering machine but a \emph{decision-making agent}.

At the same time, mobile manipulation robots have reached a level of mechanical
maturity that makes them credible platforms for performing physical tasks in real
environments.
The missing ingredient was the cognitive layer --- and large-scale VLMs now provide it
off the shelf, without the years of task-specific training data that previous approaches
required.

The natural convergence of these two trajectories is the \textbf{embodied AI agent}:
a physical robot that perceives its environment through cameras and sensors, reasons
about what to do using a VLM, and executes multi-step manipulation tasks from a single
natural-language instruction.

Consider a concrete example that motivated this thesis.
A user says: \emph{``Hello TIAGo, I am thirsty, can you do something about that?''}
The robot must understand the implicit request (find and bring a bottle), search the
room if the bottle is not immediately visible, grasp it, locate the user, and hand it
over --- all without explicit programming of any of these steps.
This is the task successfully demonstrated in this work.

This thesis addresses the challenge of building such a system on the PAL Robotics
\tiago{} mobile manipulator platform --- under the additional constraint of a
production robot environment (ROS Melodic, Python 3.6 Docker) that is
incompatible with modern deep learning frameworks, requiring a novel split-process
architecture to bridge the gap.

\section{Problem Statement}

Programming a robot to perform multi-step manipulation tasks in unstructured environments
presents the following challenges:

\begin{enumerate}
  \item \textbf{Perception:} Objects must be detected and localised in 3-D space from
        noisy RGB-D sensor data, with sufficient accuracy for the gripper to close around
        them reliably.

  \item \textbf{Reasoning:} Given a natural-language instruction, the system must
        decompose the task into a sequence of executable skills, selecting the appropriate
        action at each step based on what it currently observes.

  \item \textbf{Execution:} Physical manipulation — grasping, lifting, and handing over
        objects — must be robust to sensor noise, planning failures, and kinematic
        constraints.

  \item \textbf{Safety:} The robot must not harm itself, the environment, or nearby
        humans, and must be able to recover from unexpected failures without human
        intervention.

  \item \textbf{Deployment constraints:} The target platform runs ROS Melodic with
        Python 3.6, which is incompatible with modern deep learning frameworks.
        The system must work within these constraints without sacrificing capability.
\end{enumerate}

\section{Objectives}

The primary objective of this work is to design, implement, and evaluate a
\textbf{Vision--Language--Action (VLA) agent} for the \tiago{} robot that:

\begin{enumerate}
  \item Accepts natural-language commands via voice or text input.
  \item Perceives the environment through the robot's RGB-D camera using a
        state-of-the-art object detector.
  \item Reasons about the task using a multimodal large language model.
  \item Executes a library of physical manipulation skills, including grasping,
        searching, and handing over objects.
  \item Maintains safety through affordance checking, failure counting, and safe-pose
        recovery.
  \item Operates reliably within the Python 3.6 / ROS Melodic deployment constraint.
\end{enumerate}

Secondary objectives include:

\begin{itemize}
  \item Iterative improvement of the grasping pipeline with measurable success rates.
  \item Development of a web-based live visualiser for monitoring the agent.
  \item Evaluation of multiple YOLO model variants for the perception task.
  \item Characterisation of VLM latency and skill selection accuracy.
\end{itemize}

\section{Contributions}

The main contributions of this work are:

\begin{enumerate}
  \item \textbf{A complete VLA agent for TIAGo}, integrating Gemini 2.0 Flash, YOLO11n,
        and a modular skill library in a multi-step perceive--reason--act loop.

  \item \textbf{A robust 3-D object localisation pipeline} using RANSAC cylinder
        fitting on RGB-D point clouds with 5-frame averaging for stability.

  \item \textbf{A split-process YOLO microservice architecture} that overcomes the
        Python 3.6 Docker constraint by deploying the neural network on the host machine
        and communicating via HTTP.

  \item \textbf{An iterative grasping pipeline} (v1 to v5) with documented improvements
        at each stage, achieving reliable top-down grasping using torso-descent vertical
        approach.

  \item \textbf{A handover skill} that detects a person via YOLO and depth sensing,
        inserts a person collision object in MoveIt, and presents the grasped object at
        an ergonomic handover height.

  \item \textbf{A multi-layer safety system} including affordance gating, consecutive
        failure tracking, safe-mode entry, and graceful shutdown.
\end{enumerate}

\section{Report Structure}

The remainder of this report is organised as follows:

\begin{description}
  \item[Chapter~\ref{ch:sota}] reviews the state of the art in embodied AI, VLMs,
        object detection, and mobile manipulation.
  \item[Chapter~\ref{ch:system}] describes the complete system: hardware platform,
        overall architecture, perception pipeline, VLM reasoning, manipulation and
        grasping, and agent orchestration --- presented as a unified narrative that
        follows a single command through the system.
  \item[Chapter~\ref{ch:experiments}] presents experimental results, evaluating
        detection speed, grasping success rate, VLM latency, and end-to-end task
        completion.
  \item[Chapter~\ref{ch:conclusion}] discusses findings, limitations, and comparisons
        with related work, then concludes with contributions and future directions.
\end{description}
